\documentclass[12pt]{article}
\usepackage[utf8]{inputenc}
\usepackage[T1]{fontenc}
\usepackage[margin=1in]{geometry}
\usepackage{setspace}
\usepackage{xcolor}
\usepackage[colorlinks=true, urlcolor=blue!60!black, linkcolor=blue!60!black]{hyperref}
\setlength{\parindent}{0em}
\setlength{\parskip}{0.45em}
\renewcommand{\baselinestretch}{1.0}
\begin{document}
\begin{center}
{\Large \textit{United States v. Virginia}, 518 U.S. 515 (1996)\par}
\vspace{0.3em}
{\normalsize Case Brief --- One-Page Edition\par}
\end{center}

\textbf{Memory Jogger.} VMI must admit women: no ``exceedingly persuasive justification,'' and a separate women's program was a ``pale shadow'' (\href{https://supreme.justia.com/cases/federal/us/518/515/}{518 U.S. 515}).

\textbf{Facts.} VMI, Virginia's only single-sex public college, trained men through an ``adversative method.'' After the United States sued, Virginia created a weaker women's program (VWIL) at private Mary Baldwin College.

\textbf{Procedural History.} The Fourth Circuit found VMI's policy invalid but later approved VWIL.

\textbf{Votes.} 7--1: Ginsburg for the Court; Rehnquist concurring in the judgment; Scalia dissenting; Thomas recused.

\textbf{Holding.} Sex classifications require an exceedingly persuasive, genuine justification free of overbroad generalizations; an unequal separate program is no remedy.

\textbf{Reasoning.} VMI was not founded for diversity; fears that women would ruin its method were ``self-fulfilling prophec[ies]''; VWIL lacked VMI's resources and prestige. Rehnquist would have accepted a truly equal women's school; Scalia said tradition should yield only to democracy.

\textbf{Examples.} \emph{Same-Side:} a male-only state maritime academy. \emph{Opposite-Side:} a compensatory fellowship for women engineers. \emph{Fence-Sitter:} twin, equal single-sex high schools.

\textbf{Critique.} Critics call the test ``strict scrutiny in disguise''; admirers see an anti-hierarchy vision of equality.

\textbf{Key Quotation.} ``The justification must be genuine, not hypothesized or invented post hoc in response to litigation. And it must not rely on overbroad generalizations about the different talents, capacities, or preferences of males and females.'' --- Justice Ginsburg, 518 U.S. at 533.

\end{document}
